\subsection{The leafwise path category}

In~\cite{mac3}, we introduced a diffeological version of the Moore path category for any regular folia\-tion~$(M,\FF)$. The~objects of this category are simply points in $M$, while the morphisms are smooth, leafwise paths which have \textit{sitting instants} in that they are constant in small neighbourhoods of their endpoints. Composition of morphisms in this category is simply concatenation of paths. In~\cite{VilGar1}, the authors introduce an analogous diffeological space for \textit{singular} foliations, however concatenation of paths in this space no longer defines a category. In~this section, we~intro\-duce a hybrid of these two approaches~-- a diffeological space of integral curves of vector fields defining a singular foliation, for which concatenation of paths defines an associative and smooth multiplication.

We begin by recalling the definition of the path category of a diffeological space from~\cite{mac3} (cf.~\cite{ptfunctor}).

\begin{Definition}%\label{pathcat}
Let $\XX$ be a diffeological space. The~\textit{path category of $\XX$} is the diffeological sub\-space $\PP(\XX)$ of the diffeological product $C^{\infty}(\RB_{\geq0},\XX)\times\RB_{\geq0}$ consisting of pairs $(\gamma,d)$ for which there exist neighbourhoods of $0$ and of $[d,\infty)$ in $\RB_{\geq0}$ on which $\gamma$ is constant. Path categories are functorial~-- given any smooth map $f\colon \XX\rightarrow \YY$ of diffeological spaces, the formula
\begin{gather*}
\PP(f)(\gamma,d):=(f\circ\gamma,d)
\end{gather*}
defines a smooth functor $\PP(f)\colon \PP(\XX)\rightarrow \PP(\YY)$ of diffeological categories.
\end{Definition}

Given any diffeological space $\XX$, range and source maps $r$ and $s$ mapping $\PP(\XX)\rightarrow \XX$ are defined respectively by $(\gamma,d)\mapsto \gamma(d)$ and $(\gamma,d)\mapsto\gamma(0)$, and whenever $r(\gamma_{2},d_{2}) = s(\gamma_{1},d_{1})$, we~define the product $(\gamma_{1}\gamma_{2},d_{1}+d_{2})$ of $(\gamma_{1},d_{1})$ and $(\gamma_{2},d_{2})$ by the formula
\begin{gather*}
\gamma_{1}\gamma_{2}(t):=
\begin{cases}
\gamma_{2}(t),&\text{for}\quad 0\leq t\leq d_{2},
\\
\gamma_{1}(t-d_{2}),&\text{for}\quad d_{2}\leq t<\infty.
\end{cases}
\end{gather*}
This product, together with the range and source maps, are smooth, so that $\PP(\XX)$ is a dif\-feo\-lo\-gi\-cal category~\cite[Proposition~3.22]{mac3}. Moreover~\cite[Proposition~3.23]{mac3} there is a smooth involution $\iota\colon \PP(\XX)\ni(\gamma,d)\mapsto\big(\gamma^{-1},d\big)\rightarrow\PP(\XX)$ defined by the formula
\begin{gather*}
\gamma^{-1}(t):=
\begin{cases}
\gamma(d-t),&\text{if}\quad 0\leq t\leq d,
\\
\gamma(0),&\text{for}\quad t\geq d.
\end{cases}
\end{gather*}
Under favourable circumstances, which will be explicated in this section, the involution $\iota$ des\-cends to a genuine inversion on certain diffeological quotients of $\PP(\XX)$, giving such quotients the structures of diffeological groupoids.

Suppose in particular that $(M,\FF)$ is a singularly foliated manifold, and let $\Gamma_{\g}(\FF)$ be the pseudo-bundle of germs of $\FF$. By connectedness of $\RB_{\geq0}:=[0,\infty)$, any smooth map $\tilde{\gamma}\colon [0,\infty)\rightarrow\Gamma_{\g}(\FF)$ has the form
\begin{gather}\label{notation}
\tilde{\gamma}(t) = \big(\gamma(t),[X(t)]^{\g}_{\gamma(t)}\big),\qquad t\in[0,\infty),
\end{gather}
where $\gamma\colon [0,\infty)\rightarrow M$ is a smooth curve, and where $X\colon [0,\infty)\rightarrow\Gamma_{\loc}(\FF)$ is a smooth function. We~will implicitly use the notation of equation~\eqref{notation} in what follows.

\begin{Definition}\label{leafwisepaths}
Let $(M,\FF)$ be a singularly foliated manifold. We~define the \textit{leafwise} or \textit{$\FF$-path category} $\PP(\FF)$ to be the diffeological subspace of $\PP(\Gamma_{\g}(\FF))$ consisting of triples $\big(\gamma,[X]^{\g},d\big)$ for which $X\colon [0,\infty)\rightarrow\Gamma_{\loc}(\FF)$ satisfies
\begin{enumerate}\itemsep=0pt
\item[(1)] $\dom(X(t))$ is equal to a fixed open neighbourhood of $\gamma([0,\infty)) = \gamma([0,d])$ for all $t$,
\item[(2)] $X(t)(\gamma(t)) = \dot{\gamma}(t)$ for all $t\in[0,\infty)$, and
\item[(3)] $[X(0)]^{\g}_{\gamma(0)} = 0$ and $[X(t)]^{\g}_{\gamma(t)} = 0$ for all $t\geq d$.
\end{enumerate}
With range and source onto $M$ given by $r\big(\gamma,[X]^{\g},d\big):=\gamma(d)$ and $s\big(\gamma,[X]^{\g},d\big):=\gamma(0)$ respectively, by item 3, $\PP(\FF)$ inherits composition from $\PP(\Gamma_{\g}(\FF))$ to become a diffeological category with object space $M$.
\end{Definition}

Definition~\ref{leafwisepaths} is in practice the same as~\cite[Definition~3.1]{VilGar1} given by Garmendia--Villatoro. Note crucially that Definition~\ref{leafwisepaths} requires strictly more information than just the path in~$M$-re\-quiring in addition a time-dependent extension of the tangent field of $\gamma$ to an open neighbourhood of $\gamma$. This is so that flows of elements of $\PP(\FF)$ determine germs of diffeomorphisms defined in open neighbourhoods of their sources. Thus Definition~\ref{leafwisepaths} may be contrasted with the simpler~\cite[Definition~3.23]{mac3} for regular foliations, where such an extension is not explicitly required. In~the regular case, the tangent field along a leafwise path can always be \textit{canonically} extended to a tangent field in a (transverse) open neighbourhood of the path in any foliated chart. The~same is not true in the singular setting.
